\newcommand{\TomTat}{Đồ án ước lượng trạng thái sức khoẻ (SOH) của cell lithium-ion từ 16 đặc trưng đoạn sạc khi chỉ một phần nhỏ cell được đo dung lượng, trên bốn bộ dữ liệu công khai (387 cell). Phân tích đặc trưng bằng tương quan hạng, bỏ từng nhóm đầu vào, độ quan trọng hoán vị và chọn trên tập xác thực cho thấy cả ba nhóm đặc trưng sạc và chỉ số chu kỳ đều mang thông tin. Chuẩn hoá theo cả vòng đời cell và lọc ngoại lai bằng tương lai của cell bị loại vì là rò rỉ. Kiểm toán cho thấy chia theo chu kỳ làm sai số thấp đi 1{,}2--1{,}8 lần và chuẩn hoá theo cả vòng đời cell làm sai số đổi 0{,}46--2{,}76 lần tuỳ bộ. Trong giao thức không rò rỉ (kiểm định chéo 5 fold theo cell, lặp lại, kiểm định ghép cặp), PINN bán giám sát có sai số bằng 0{,}39--0{,}90 lần MLP ở 10\,\% nhãn, nhưng GBDT vẫn tốt hơn nó khi có từ 30\,\% nhãn. Mô hình kết hợp PINN-lai, cấu hình chốt trước khi đọc tập kiểm tra, có sai số bằng 0{,}67--0{,}95 lần PINN-semi và thấp hơn GBDT ở 10/16 trường hợp. Trên đúng 39 nhóm đánh giá và với tiền xử lý của bài PINN4SOH, PINN-lai thấp hơn số công bố ở 11/12 chế độ và thấp hơn mã gốc của bài chạy lại ở 12/12, nhưng một GBDT cùng tiền xử lý đạt kết quả tương tự. Với tiền xử lý không rò rỉ, PINN-lai thấp hơn số công bố trên TJU và MIT, cao hơn trên HUST. Đo thưa mọi cell với 1{,}2--6{,}0\,\% số chu kỳ cho sai số bằng 0{,}85--0{,}93 lần gắn nhãn đầy đủ cho 30\,\% số cell. Đưa lịch sử của cell vào dự báo, đồ án xây dựng HiGBDT (GBDT có đặc trưng lịch sử SOH) và HiPINN, mô hình lấy trung vị của HiGBDT, PINN-dự báo và một nhánh theo đời cell. Khi có lịch sử đo dung lượng, HiPINN có sai số dự báo SOH bằng 0{,}62--0{,}91 lần PINN-GBDT, mô hình dự báo không có đặc trưng lịch sử.}
